% !TEX program = lualatex
% Transparencias de Fundamentos de las Matemáticas II
% Clase 12: Números racionales

\documentclass[aspectratio=169,11pt]{beamer}

\usepackage{fontspec}
\usepackage[spanish,es-nodecimaldot]{babel}
\usepackage{amsmath,amssymb,mathtools}
\usepackage{booktabs}
\usepackage{csquotes}
\usepackage{xcolor}

\usetheme{Madrid}
\usecolortheme{default}
\setbeamertemplate{navigation symbols}{}
\setbeamertemplate{footline}[frame number]

\definecolor{FdMBlue}{RGB}{20,55,100}
\definecolor{FdMRed}{RGB}{130,30,30}
\definecolor{FdMGreen}{RGB}{25,100,60}

\setbeamercolor{structure}{fg=FdMBlue}
\setbeamercolor{title}{fg=white,bg=FdMBlue}
\setbeamercolor{frametitle}{fg=white,bg=FdMBlue}
\setbeamercolor{block title}{fg=white,bg=FdMBlue}
\setbeamercolor{block body}{bg=blue!3}
\setbeamercolor{alerted text}{fg=FdMRed}

\setmainfont{Latin Modern Roman}
\setsansfont{Latin Modern Sans}
\setmonofont{Latin Modern Mono}

\newcommand{\N}{\mathbb{N}}
\newcommand{\Z}{\mathbb{Z}}
\newcommand{\Q}{\mathbb{Q}}
\newcommand{\R}{\mathbb{R}}
\newcommand{\C}{\mathbb{C}}
\newcommand{\Pow}{\mathcal{P}}
\newcommand{\id}{\operatorname{id}}
\newcommand{\mcd}{\operatorname{mcd}}
\newcommand{\mcm}{\operatorname{mcm}}
\newcommand{\im}{\operatorname{Im}}

\title[FdM II -- Clase 12]{Números racionales}
\subtitle{Fracciones como clases de equivalencia}
\author{Antonio Falcó Montesinos}
\institute{Universidad CEU Cardenal Herrera}
\date{Fundamentos de las Matemáticas II}

\begin{document}

\begin{frame}
  \titlepage
\end{frame}

\begin{frame}{Objetivos de la clase}
\begin{itemize}
  \item Motivar la construcción de \(\mathbb{Q}\).
  \item Entender fracciones equivalentes.
  \item Definir suma, producto, orden y densidad.
\end{itemize}
\end{frame}

\begin{frame}{Mapa de la clase}
\tableofcontents
\end{frame}

\section{Construcción}

\begin{frame}{Construcción}
\begin{block}{Motivación}
\begin{itemize}
  \item En \(\mathbb{Z}\), no siempre se puede dividir.
  \item La ecuación \(2x=1\) no tiene solución entera.
  \item Los racionales amplían \(\mathbb{Z}\) para permitir cocientes.
\end{itemize}
\end{block}
\end{frame}

\begin{frame}{Construcción}
\begin{block}{Pares admisibles}
\begin{itemize}
  \item Se consideran pares \((a,b)\in\mathbb{Z}\times\mathbb{Z}\) con \(b\neq 0\).
  \item El par representa informalmente la fracción \(a/b\).
  \item Distintos pares pueden representar el mismo racional.
\end{itemize}
\end{block}
\end{frame}

\begin{frame}{Construcción}
\begin{block}{Equivalencia}
\begin{itemize}
  \item \((a,b)\sim(c,d)\) si \(ad=bc\).
  \item Un racional es una clase de fracciones equivalentes.
  \item \(\frac{1}{2}\), \(\frac{2}{4}\) y \(\frac{-3}{-6}\) representan la misma clase.
\end{itemize}
\end{block}
\end{frame}

\section{Estructura}

\begin{frame}{Estructura}
\begin{block}{Operaciones}
\begin{itemize}
  \item \(\frac{a}{b}+\frac{c}{d}=\frac{ad+bc}{bd}\).
  \item \(\frac{a}{b}\frac{c}{d}=\frac{ac}{bd}\).
  \item Hay que comprobar que estas operaciones están bien definidas sobre clases.
\end{itemize}
\end{block}
\end{frame}

\begin{frame}{Estructura}
\begin{block}{Orden y densidad}
\begin{itemize}
  \item El orden puede definirse comparando diferencias positivas.
  \item Entre dos racionales distintos siempre hay otro racional.
  \item Si \(r<s\), entonces \(q=(r+s)/2\) cumple \(r<q<s\).
\end{itemize}
\end{block}
\end{frame}

\section{Profundización}

\begin{frame}{Profundización}
\begin{block}{Enteros dentro de racionales}
\begin{itemize}
  \item Cada \(n\in\mathbb{Z}\) se identifica con \(n/1\).
  \item Así, \(\mathbb{Z}\subset\mathbb{Q}\) en sentido estructural.
  \item Esta identificación conserva suma y producto.
\end{itemize}
\end{block}
\end{frame}

\begin{frame}{Profundización}
\begin{block}{Decimales}
\begin{itemize}
  \item Los racionales tienen expansiones decimales finitas o periódicas.
  \item Un decimal no periódico anticipa números irracionales.
  \item La representación decimal no debe reemplazar la definición como clase.
\end{itemize}
\end{block}
\end{frame}

\begin{frame}{Profundización}
\begin{block}{Error frecuente}
\begin{itemize}
  \item Pensar que una fracción es solo un par \((a,b)\).
  \item Olvidar que el denominador debe ser no nulo.
  \item Comparar productos cruzados sin controlar signos de denominadores.
\end{itemize}
\end{block}
\end{frame}


\section{Detalle y práctica}

\begin{frame}{Detalle y práctica}
\begin{block}{Fracciones como clases}
\begin{itemize}
  \item En \(\mathbb{Q}\), las fracciones \(\frac{a}{b}\) y \(\frac{c}{d}\) representan el mismo número si \(ad=bc\).
  \item La condición \(b,d\neq 0\) es esencial.
  \item Reducir una fracción no cambia el racional, solo cambia su representante.
\end{itemize}
\end{block}
\end{frame}

\begin{frame}{Detalle y práctica}
\begin{block}{Actividad breve}
\begin{itemize}
  \item Comprobar si \(\frac{18}{24}=\frac{3}{4}\).
  \item Ordenar \(-\frac{2}{3}\), \(\frac{1}{2}\) y \(-\frac{5}{6}\).
  \item Justificar por qué todo entero \(z\) se identifica con \(\frac{z}{1}\).
\end{itemize}
\end{block}
\end{frame}

\begin{frame}{Cierre}
\begin{itemize}
  \item Los racionales son clases de fracciones equivalentes.
  \item \(\mathbb{Q}\) permite dividir por no nulos.
  \item La densidad no significa completitud.
\end{itemize}
\end{frame}

\begin{frame}{Trabajo recomendado}
\begin{itemize}
  \item Releer las secciones correspondientes del manual.
  \item Identificar definiciones, símbolos nuevos y enunciados principales.
  \item Preparar dudas y ejemplos para el seminario asociado.
\end{itemize}
\end{frame}

\end{document}
